% ---------- פרק 6: תוכנית עבודה ----------
% לפי פרק 8 במדריך. הספירה מתחילה ב-11.11.2026; אבן דרך 1 = הגשת הדוח המכין ב-25.11.2026 (עמית, 16.11).
% RACI לפי עמית (16.11): מנחה תמיד A; תע"א C במשימות 1,3,6,8 ו-I במשימות 7,9. 22.11: עמודות בשמות (עמית סוסן, טל אזואלוס); R/R ב-T0, T1.3, T5.1 נשאר לבקשת עמית.
\chapter{תוכנית עבודה}\label{ch:workplan}

הפרק מתאר איך המתודולוגיה תמומש, ועונה על שלוש שאלות: מה עושים, מי עושה ומתי. לוח הזמנים מתחיל ב-11.11.2026 ומתפרס על 16 שבועות, עם שלוש אבני דרך:
\begin{enumerate}
\item \textbf{הגשת הדוח המכין} ב-25.11.2026, בסוף שבוע 2.
\item \textbf{אב-טיפוס עובד} בסוף שבוע 10, לפי יעד ה-\en{SMART} בפרק~\ref{ch:problem}.
\item \textbf{השלמת התיקוף} בסוף שבוע 13.
\end{enumerate}

\section{פירוק העבודה (\en{WBS})}\label{sec:wbs}

הפרויקט מחולק לחמש חבילות עבודה (\en{Work Packages}). לכל חבילה יש משימות ותוצר מוגדר (איור~\ref{fig:wbs} וטבלה~\ref{tab:wbs}).

\begin{figure}[!htbp]
\centering
\begin{tikzpicture}[
  wp/.style={draw, rounded corners, fill=blue!6, align=center, font=\footnotesize, text width=2.4cm, minimum height=12mm},
  top/.style={draw, rounded corners, fill=gray!15, align=center, font=\small\bfseries, text width=4.5cm, minimum height=9mm}]
\node[top] (p) at (0,0) {\foreignlanguage{hebrew}{אב-טיפוס מצפן ספקים}};
\foreach \i/\t [count=\k] in {
  5/{תיעוד והגנה},
  4/{תיקוף},
  3/{שקלול ודוח},
  2/{שער וסוכנים},
  1/{נתונים והכנה}}{
  \node[wp] (w\i) at ({(\k-3)*2.8},-2) {\en{WP\i}\\\foreignlanguage{hebrew}{\t}};
  \draw[thick] (p.south) -- ++(0,-0.4) -| (w\i.north);
}
\end{tikzpicture}
\caption{מבנה פירוק העבודה: חמש חבילות תחת פרויקט אחד}\label{fig:wbs}
\end{figure}

{\small
\begin{longtable}{|>{\raggedright\arraybackslash}p{1.3cm}|>{\raggedright\arraybackslash}p{2.6cm}|>{\raggedright\arraybackslash}p{5.6cm}|>{\raggedright\arraybackslash}p{4.3cm}|}
\caption{חבילות העבודה, המשימות והתוצרים}\label{tab:wbs}\\
\hline
\textbf{מזהה} & \textbf{החבילה} & \textbf{משימות} & \textbf{תוצר} \\ \hline
\endfirsthead
\hline
\textbf{מזהה} & \textbf{החבילה} & \textbf{משימות} & \textbf{תוצר} \\ \hline
\endhead
\en{WP1} & נתונים והכנה & \en{T1.1} בניית מאגר התיקוף: ספקים שכשלו וספקים מוכחים, עם תוצאה צפויה לכל מקרה (טבלה~\ref{tab:cases}). \en{T1.2} חיבור למקורות המידע ובדיקת תנאי השימוש שלהם. \en{T1.3} מילוי \en{AHP} וניתוח רגישות & מאגר תיקוף מתועד, מחברים למקורות, משקלות סופיים \\ \hline
\en{WP2} & שער וסוכנים & \en{T2.1} זיהוי הישות ושער הסנקציות. \en{T2.2} שבעה סוכני ממד: ארבעה לפרופיל החוזקות ושלושה לדגלי הסיכון, עם פלט במבנה קבוע & קוד מתועד ובדיקות יחידה \\ \hline
\en{WP3} & שקלול ודוח & \en{T3.1} מנוע השקלול, כללי ההחלטה וכותב הדוח & אב-טיפוס שמפיק דוח מלא \\ \hline
\en{WP4} & תיקוף & \en{T4.1} בדיקה לאחור לפי נקודת זמן על ספקים שכשלו ועל ספקים מוכחים, עם שני מודלי שפה לפחות, וכיול הספים על קבוצת הכיול. \en{T4.2} שאלון משוב מהתעשייה האווירית (\en{H3}) & דוח תיקוף: \en{H1} עד \en{H3} \\ \hline
\en{WP5} & תיעוד והגנה & \en{T5.1} ספר הפרויקט, מצגת והגנה & הגשה סופית \\ \hline
\end{longtable}
}

\section{מטריצת אחריות (\en{RACI})}\label{sec:raci}

טבלה~\ref{tab:raci} מגדירה לכל משימה את התפקיד של כל משתתף: \textbf{\en{R}} (\en{Responsible}) מבצעת את המשימה ואחראית לתוצר; \textbf{\en{A}} (\en{Accountable}) מאשר את התוצר; \textbf{\en{C}} (\en{Consulted}) מתייעצים איתו במהלך הביצוע; \textbf{\en{I}} (\en{Informed}) מקבל עדכון על התוצאה.

{\small
\begin{longtable}{|c|>{\raggedright\arraybackslash}p{5.6cm}|c|c|c|c|}
\caption{מטריצת אחריות}\label{tab:raci}\\
\hline
\textbf{\#} & \textbf{משימה} & \textbf{עמית סוסן} & \textbf{טל אזואלוס} & \textbf{מנחה אקדמי} & \textbf{תע"א} \\ \hline
\endfirsthead
\hline
\textbf{\#} & \textbf{משימה} & \textbf{עמית סוסן} & \textbf{טל אזואלוס} & \textbf{מנחה אקדמי} & \textbf{תע"א} \\ \hline
\endhead
0 & \en{T0} כתיבת הדוח המכין & \en{R} & \en{R} & \en{A} & --- \\ \hline
1 & \en{T1.1} מאגר תיקוף: ספקים שכשלו וספקים מוכחים & \en{R} & \en{C} & \en{A} & \en{C} \\ \hline
2 & \en{T1.2} מחברים למקורות מידע & \en{C} & \en{R} & \en{A} & --- \\ \hline
3 & \en{T1.3} \en{AHP} וניתוח רגישות & \en{R} & \en{R} & \en{A} & \en{C} \\ \hline
4 & \en{T2.1} זיהוי ישות ושער סנקציות & \en{C} & \en{R} & \en{A} & --- \\ \hline
5 & \en{T2.2} סוכני הממדים: חוזקות וסיכונים & \en{R} & \en{C} & \en{A} & --- \\ \hline
6 & \en{T3.1} שקלול ודוח & \en{C} & \en{R} & \en{A} & \en{C} \\ \hline
7 & \en{T4.1} בדיקה לאחור וכיול & \en{R} & \en{C} & \en{A} & \en{I} \\ \hline
8 & \en{T4.2} שאלון משוב & \en{C} & \en{R} & \en{A} & \en{C} \\ \hline
9 & \en{T5.1} ספר, מצגת והגנה & \en{R} & \en{R} & \en{A} & \en{I} \\ \hline
\end{longtable}
}

\section{לוח זמנים ותרשים גאנט}\label{sec:gantt}

טבלה~\ref{tab:tasks} מפרטת את המשימות, את התלויות ביניהן ואת אבני הדרך. איור~\ref{fig:gantt} מציג אותן בתרשים גאנט. שבוע 1 מתחיל ב-11.11.2026.

\textbf{הנתיב הקריטי} הוא שרשרת המשימות שבה כל עיכוב מעכב את כל הפרויקט:
\begin{center}
\en{T0 $\rightarrow$ T1.2 $\rightarrow$ T2.1 $\rightarrow$ T2.2 $\rightarrow$ T3.1 $\rightarrow$ T4.1 $\rightarrow$ T4.2 $\rightarrow$ \foreignlanguage{hebrew}{רזרבה} $\rightarrow$ T5.1}
\end{center}
אורך הנתיב הוא \en{$2+3+1+3+1+2+1+1+2=16$} שבועות.

ל-\en{T1.1} ול-\en{T1.3} יש מרווח: הן צריכות להסתיים רק לפני \en{T4.1} ו-\en{T3.1}. כך, עיכוב קצר בבניית מאגר התיקוף לא יעכב את הפרויקט.

\begin{figure}[!htbp]
\centering
\begin{tikzpicture}[x=0.66cm, y=0.46cm]
% ציר שבועות
\foreach \w in {1,...,16}{
  \node[font=\scriptsize] at (\w-0.5,0.7) {\en{\w}};
  \draw[gray!25] (\w,0.3) -- (\w,-11.3);
}
\draw[gray!25] (0,0.3) -- (0,-11.3);
% משימות: שורה, שם, התחלה, סוף, קריטית?
\foreach \row/\name/\s/\e/\c in {
  1/T0/1/2/1,
  2/T1.1/3/4/0,
  3/T1.2/3/5/1,
  4/T1.3/4/5/0,
  5/T2.1/6/6/1,
  6/T2.2/7/9/1,
  7/T3.1/10/10/1,
  8/T4.1/11/12/1,
  9/T4.2/13/13/1,
  10/{\foreignlanguage{hebrew}{רזרבה}}/14/14/1,
  11/T5.1/15/16/1}{
  \node[font=\scriptsize, anchor=east] at (-0.2,-\row+0.5) {\name};
  \ifnum\c=1
    \fill[red!65] (\s-1,-\row+0.15) rectangle (\e,-\row+0.85);
  \else
    \fill[blue!45] (\s-1,-\row+0.15) rectangle (\e,-\row+0.85);
  \fi
}
% אבני דרך
\foreach \x/\y/\n in {2/-0.5/1, 10/-6.5/2, 13/-8.5/3}{
  \fill[green!50!black] (\x,\y) +(0,0.3) -- +(0.3,0) -- +(0,-0.3) -- +(-0.3,0) -- cycle;
  \node[font=\scriptsize, anchor=west] at (\x+0.3,\y) {\foreignlanguage{hebrew}{אבן דרך \n}};
}
% מקרא
\fill[red!65] (0,-12.3) rectangle (0.6,-11.9); \node[font=\scriptsize, anchor=west] at (0.7,-12.1) {\foreignlanguage{hebrew}{נתיב קריטי}};
\fill[blue!45] (5,-12.3) rectangle (5.6,-11.9); \node[font=\scriptsize, anchor=west] at (5.7,-12.1) {\foreignlanguage{hebrew}{משימה עם מרווח}};
\end{tikzpicture}
\caption{תרשים גאנט ל-16 שבועות, החל מ-11.11.2026. הנתיב הקריטי באדום, ושלוש אבני הדרך בירוק.}\label{fig:gantt}
\end{figure}

{\small
\begin{longtable}{|l|>{\raggedright\arraybackslash}p{5cm}|c|>{\raggedright\arraybackslash}p{2.2cm}|>{\raggedright\arraybackslash}p{3.4cm}|}
\caption{משימות, תלויות ואבני דרך}\label{tab:tasks}\\
\hline
\textbf{קוד} & \textbf{המשימה} & \textbf{שבועות} & \textbf{קודמות} & \textbf{אבן דרך / תוצר} \\ \hline
\endfirsthead
\hline
\textbf{קוד} & \textbf{המשימה} & \textbf{שבועות} & \textbf{קודמות} & \textbf{אבן דרך / תוצר} \\ \hline
\endhead
\en{T0} & כתיבת הדוח המכין & \en{1--2} & --- & \textbf{1: הגשת הדוח המכין (25.11.2026)} \\ \hline
\en{T1.1} & מאגר תיקוף: ספקים שכשלו וספקים מוכחים & \en{3--4} & \en{T0} & מאגר תיקוף \\ \hline
\en{T1.2} & מחברים למקורות מידע ובדיקת תנאי שימוש & \en{3--5} & \en{T0} & גישה לכל מקור \\ \hline
\en{T1.3} & \en{AHP} וניתוח רגישות & \en{4--5} & \en{T0} & משקלות סופיים \\ \hline
\en{T2.1} & זיהוי ישות ושער סנקציות & 6 & \en{T1.2} & שער עובד \\ \hline
\en{T2.2} & שבעה סוכני ממד & \en{7--9} & \en{T2.1} & פלט לכל ממד \\ \hline
\en{T3.1} & שקלול, כללי החלטה ודוח & 10 & \en{T2.2}, \en{T1.3} & \textbf{2: אב-טיפוס עובד} \\ \hline
\en{T4.1} & בדיקה לאחור בשני מודלים וכיול ספים & \en{11--12} & \en{T3.1}, \en{T1.1} & תוצאות \en{H1}, \en{H2} \\ \hline
\en{T4.2} & שאלון משוב (\en{H3}) & 13 & \en{T4.1} & \textbf{3: התיקוף הושלם} \\ \hline
--- & רזרבה & 14 & --- & --- \\ \hline
\en{T5.1} & ספר פרויקט, מצגת והגנה & \en{15--16} & \en{T4.2} & הגשה והגנה \\ \hline
\end{longtable}
}

\textbf{מעקב.} כל פגישה עם המנחה מתועדת בפרוטוקול קצר: מה הושלם, מה התעכב ומה הצעד הבא. הקוד והדוח מנוהלים במאגר \en{GitHub}, וכך אפשר לראות בכל רגע את מצב העבודה מול התוכנית.
